\chapter{Circuitos en serie}\label{ch:g4:series-circuits}

El año pasado, una guirnalda de luces de fiesta se quedó a oscuras
porque faltaba una sola \omterm{def:g3:first-electric-circuit:bulb}{bombilla} --- y prometimos que su colocación
tenía nombre. Aquí está: el \omterm{def:g4:series-circuits:series}{circuito en serie}, la manera más sencilla
de poner muchos jugadores en un solo bucle. Tiene reglas de la casa
estrictas, y una linterna de tu cajón lleva toda la vida
obedeciéndolas.

\section{Un bucle, varios jugadores}

\begin{definition}[En serie]\label{def:g4:series-circuits:series}
Los jugadores de un circuito están \emph{en serie}\index{circuito en
serie} cuando se sientan en un único bucle, uno detrás de otro, como
las cuentas de un collar: la carretera de la corriente pasa por el
primero, después por el segundo y después por el tercero, sin
bifurcarse en ningún sitio. Un circuito construido así es un
\emph{circuito en serie}.
\end{definition}

\begin{omfigure}
\begin{tikzpicture}[scale=0.9]
  \draw (0,0) to[battery1, l=pila] (0,2.6)
    to[short] (1.6,2.6)
    to[cspst, l=interruptor] (3.2,2.6)
    to[lamp, l=bombilla 1] (5.2,2.6)
    to[lamp, l=bombilla 2] (7.2,2.6)
    to[short] (7.2,0)
    to[short] (0,0);
\end{tikzpicture}

{\small Un \omterm{def:g4:series-circuits:series}{circuito en serie}: \omterm{def:g3:first-electric-circuit:battery}{pila}, \omterm{ex:g3:first-electric-circuit:switch}{interruptor} y dos
\omterm{def:g3:first-electric-circuit:bulb}{bombillas} en un único bucle --- todos los jugadores ensartados en el
mismo collar.}
\end{omfigure}

\begin{method}[Construirlo]\label{met:g4:series-circuits:build}
Con una \omterm{def:g3:first-electric-circuit:battery}{pila}, dos \omterm{def:g3:first-electric-circuit:bulb}{bombillas}, un \omterm{ex:g3:first-electric-circuit:switch}{interruptor} y cables:
\begin{enumerate}
  \item sal del \omterm{def:g3:first-electric-circuit:battery}{borne} $+$ de la \omterm{def:g3:first-electric-circuit:battery}{pila} con un cable hasta el
        \omterm{ex:g3:first-electric-circuit:switch}{interruptor};
  \item lleva un cable del \omterm{ex:g3:first-electric-circuit:switch}{interruptor} a la primera \omterm{def:g3:first-electric-circuit:bulb}{bombilla} y otro
        de la primera \omterm{def:g3:first-electric-circuit:bulb}{bombilla} a la segunda;
  \item cierra el bucle: de la segunda \omterm{def:g3:first-electric-circuit:bulb}{bombilla} de vuelta al \omterm{def:g3:first-electric-circuit:battery}{borne}
        $-$ de la \omterm{def:g3:first-electric-circuit:battery}{pila};
  \item pulsa el \omterm{ex:g3:first-electric-circuit:switch}{interruptor} y mira las dos \omterm{def:g3:first-electric-circuit:bulb}{bombillas} a la vez.
\end{enumerate}
Recorre el bucle con el dedo antes de encender: si tu dedo tiene que
elegir en algún momento entre dos carreteras, no es un
\omterm{def:g4:series-circuits:series}{circuito en serie}.
\end{method}

\section{Las reglas de la casa}

\begin{proposition}[Reglas del bucle único]\label{prop:g4:series-circuits:rules}
En un \omterm{def:g4:series-circuits:series}{circuito en serie}:
\begin{enumerate}
  \item \emph{todo o nada}: un hueco en cualquier sitio --- un
        \omterm{ex:g3:first-electric-circuit:switch}{interruptor} abierto, una \omterm{def:g3:first-electric-circuit:bulb}{bombilla} fundida, un cable
        suelto --- y todos los jugadores paran a la vez;
  \item \emph{el orden no importa}: intercambia las \omterm{def:g3:first-electric-circuit:bulb}{bombillas}, mueve
        el \omterm{ex:g3:first-electric-circuit:switch}{interruptor} al otro lado del bucle --- todo se comporta
        exactamente igual que antes;
  \item \emph{reparto}: cuantas más \omterm{def:g3:first-electric-circuit:bulb}{bombillas} hay en el bucle, más
        débil brilla cada una; el empuje de la \omterm{def:g3:first-electric-circuit:battery}{pila} se reparte a lo
        largo de todo el collar.
\end{enumerate}
\end{proposition}

\begin{example}[Ver las reglas]\label{ex:g4:series-circuits:seeing}
Construye el circuito de dos \omterm{def:g3:first-electric-circuit:bulb}{bombillas} y pon a prueba cada regla.
Desenrosca \emph{cualquiera} de las dos \omterm{def:g3:first-electric-circuit:bulb}{bombillas}: las dos se apagan
--- regla uno (el casquillo vacío es un hueco). Intercambia las dos
\omterm{def:g3:first-electric-circuit:bulb}{bombillas}, o mueve el \omterm{ex:g3:first-electric-circuit:switch}{interruptor} entre ellas: ningún cambio ---
regla dos. Vuelve a montarlo ahora con una sola \omterm{def:g3:first-electric-circuit:bulb}{bombilla}: brilla
claramente más sola que cuando tenía compañera --- la regla tres, al
revés.
\end{example}

\begin{example}[Las luces de fiesta, resueltas]\label{ex:g4:series-circuits:partylights}
El viejo misterio ya no es ningún misterio: las \omterm{def:g3:first-electric-circuit:bulb}{bombillas} de fiesta
están \omterm{def:g4:series-circuits:series}{en serie}, decenas de cuentas en un solo collar. Una
\omterm{def:g3:first-electric-circuit:bulb}{bombilla} que falta es un hueco, y la regla uno apaga la guirnalda
entera. Buscar a la \omterm{def:g3:first-electric-circuit:bulb}{bombilla} culpable probándolas una a una le
enseñó la regla uno por las malas a generaciones enteras de familias
--- las guirnaldas modernas están construidas de otra manera, como
enseñará el año que viene.
\end{example}

\begin{remark}[El interruptor manda sobre todos]\label{rem:g4:series-circuits:switch}
La regla dos tiene una consecuencia muy cómoda: un solo
\omterm{ex:g3:first-electric-circuit:switch}{interruptor}, colocado en cualquier punto del bucle, manda sobre
todos los jugadores. Por eso un solo \omterm{ex:g3:first-electric-circuit:switch}{interruptor} de pared puede
controlar toda una fila de luces del techo --- y por eso mismo no
puedes apagar solo \emph{una} \omterm{def:g3:first-electric-circuit:bulb}{bombilla} de una guirnalda \omterm{def:g4:series-circuits:series}{en serie} sin
apagarlas todas.
\end{remark}

\section{Más pilas en el collar}

\begin{example}[Nariz contra cola]\label{ex:g4:series-circuits:batteries}
Las \omterm{def:g3:first-electric-circuit:battery}{pilas} también pueden hacer cola \omterm{def:g4:series-circuits:series}{en serie} --- pero son más
quisquillosas que las \omterm{def:g3:first-electric-circuit:bulb}{bombillas}: tienen que colocarse \emph{nariz
contra cola}, con el $+$ de cada \omterm{def:g3:first-electric-circuit:battery}{pila} tocando el $-$ de la
siguiente. Dos \omterm{def:g3:first-electric-circuit:battery}{pilas} en cola de esta manera empujan juntas, y la
\omterm{def:g3:first-electric-circuit:bulb}{bombilla} brilla claramente más que con una sola. Da la vuelta a una
\omterm{def:g3:first-electric-circuit:battery}{pila} --- nariz contra nariz --- y los dos empujes se pelean entre
sí: la \omterm{def:g3:first-electric-circuit:bulb}{bombilla} no da nada.
\end{example}

\begin{omfigure}
\begin{tikzpicture}[scale=0.8]
  \draw (0,0) to[battery1] (0,1.8)
    to[battery1, l=dos pilas] (0,3.6)
    to[short] (3.6,3.6)
    to[lamp, l=¡más brillante!] (3.6,0)
    to[short] (0,0);
  \begin{scope}[xshift=6.9cm]
    \draw (0,0) to[battery1, l=una pila] (0,1.8)
      to[short] (0,3.6)
      to[short] (3.6,3.6)
      to[lamp, l=más floja] (3.6,0)
      to[short] (0,0);
  \end{scope}
\end{tikzpicture}

{\small Dos \omterm{def:g3:first-electric-circuit:battery}{pilas} en cola, nariz contra cola, empujan juntas: la misma
\omterm{def:g3:first-electric-circuit:bulb}{bombilla} brilla más que con una \omterm{def:g3:first-electric-circuit:battery}{pila} sola.}
\end{omfigure}

\begin{example}[Por dentro de la linterna]\label{ex:g4:series-circuits:flashlight}
Abre una linterna: las \omterm{def:g3:first-electric-circuit:battery}{pilas} entran una detrás de otra en un tubo ---
una cola nariz contra cola ya montada ---, y después la \omterm{def:g3:first-electric-circuit:bulb}{bombilla} y el
pulsador completan un único bucle a través de la carcasa. Una
linterna \emph{es} un \omterm{def:g4:series-circuits:series}{circuito en serie} de bolsillo. Y ahora ya sabes
por qué se queda apagada si una sola \omterm{def:g3:first-electric-circuit:battery}{pila} se mete del revés: nariz
contra nariz, los empujes se anulan.
\end{example}

\section{Ejercicios}

\begin{exercise}[$\star$]\label{exo:g4:series-circuits:1}
¿Qué significa ``\omterm{def:g4:series-circuits:series}{en serie}''? ¿Qué imagen de la vida diaria usa el
capítulo para explicarlo?
\end{exercise}

\begin{exercise}[$\star$]\label{exo:g4:series-circuits:2}
Enuncia las tres reglas de la casa del \omterm{def:g4:series-circuits:series}{circuito en serie}.
\end{exercise}

\begin{exercise}[$\star$]\label{exo:g4:series-circuits:3}
En el circuito de dos \omterm{def:g3:first-electric-circuit:bulb}{bombillas} se desenrosca la \omterm{def:g3:first-electric-circuit:bulb}{bombilla} 2. ¿Qué
le pasa a la \omterm{def:g3:first-electric-circuit:bulb}{bombilla} 1 y por qué? ¿Qué regla está actuando?
\end{exercise}

\begin{exercise}[$\star$]\label{exo:g4:series-circuits:4}
Zoe mueve el \omterm{ex:g3:first-electric-circuit:switch}{interruptor} de delante de la \omterm{def:g3:first-electric-circuit:bulb}{bombilla} 1 a detrás de
la \omterm{def:g3:first-electric-circuit:bulb}{bombilla} 2. ¿Qué cambia en el comportamiento del circuito? ¿Qué
regla lo dice?
\end{exercise}

\begin{exercise}[$\star$]\label{exo:g4:series-circuits:5}
Una \omterm{def:g3:first-electric-circuit:battery}{pila} enciende bien una \omterm{def:g3:first-electric-circuit:bulb}{bombilla}. Predice el brillo con:
dos \omterm{def:g3:first-electric-circuit:bulb}{bombillas} \omterm{def:g4:series-circuits:series}{en serie}; \ tres \omterm{def:g3:first-electric-circuit:bulb}{bombillas} \omterm{def:g4:series-circuits:series}{en serie}. ¿Qué es lo
que se está repartiendo?
\end{exercise}

\begin{exercise}[$\star$]\label{exo:g4:series-circuits:6}
¿Cómo tienen que colocarse dos \omterm{def:g3:first-electric-circuit:battery}{pilas} para empujar juntas? Dibuja la
cola con el $+$ y el $-$ marcados en cada \omterm{def:g3:first-electric-circuit:battery}{pila}.
\end{exercise}

\begin{exercise}[$\star$]\label{exo:g4:series-circuits:7}
Recorre el bucle de una linterna: nombra los jugadores que visita la
corriente desde el $+$ de la \omterm{def:g3:first-electric-circuit:battery}{pila} hasta su $-$.
\end{exercise}

\begin{exercise}[$\star$]\label{exo:g4:series-circuits:8}
El circuito de un timbre tiene una \omterm{def:g3:first-electric-circuit:battery}{pila}, un pulsador y el timbre
\omterm{def:g4:series-circuits:series}{en serie}. ¿Dónde podrías añadir un segundo pulsador para que el
timbre suene solo cuando se aprieten \emph{los dos}?
\end{exercise}

\begin{exercise}[$\star$]\label{exo:g4:series-circuits:9}
Una guirnalda de diez luces de fiesta \omterm{def:g4:series-circuits:series}{en serie} está apagada. Todas
las \omterm{def:g3:first-electric-circuit:bulb}{bombillas} parecen estar bien. Describe un plan paciente, con una
\omterm{def:g3:first-electric-circuit:bulb}{bombilla} que sepas que funciona, para encontrar la fundida.
\end{exercise}

\begin{exercise}[$\star\star$]\label{exo:g4:series-circuits:10}
La linterna de Tom lleva dos \omterm{def:g3:first-electric-circuit:battery}{pilas} y alumbra estupendamente. Después
de cambiar las \omterm{def:g3:first-electric-circuit:battery}{pilas} a oscuras, no da nada --- y sin embargo las dos
\omterm{def:g3:first-electric-circuit:battery}{pilas} nuevas están llenas, la \omterm{def:g3:first-electric-circuit:bulb}{bombilla} está bien y no hay ningún
cable suelto. ¿Qué ha hecho Tom y por qué se queda apagada la
linterna? (Mira el \cref{ex:g4:series-circuits:batteries}.)
\end{exercise}

\begin{exercise}[$\star\star$]\label{exo:g4:series-circuits:11}
Una ingeniera tiene que colocar un \omterm{ex:g3:first-electric-circuit:switch}{interruptor} de parada de
emergencia para una máquina cuyo motor está en un único bucle con su
\omterm{def:g3:first-electric-circuit:battery}{pila}. Dice: ``Sirve cualquier punto del bucle''. ¿Tiene razón? ¿Y
por qué fallaría su plan si el circuito de la máquina tuviera
carreteras con bifurcaciones? (El capítulo del año que viene les pone
nombre.)
\end{exercise}
